\documentclass[11pt]{article}
\usepackage[a4paper, margin=1in]{geometry}
\usepackage{amsmath}
\usepackage{amssymb}
\usepackage{enumitem}
\usepackage{fancyvrb}
\usepackage{xcolor}
\usepackage{hyperref}
\usepackage{booktabs}
\usepackage{array}
\usepackage{longtable}
\usepackage{fancyhdr}

\setlength{\parindent}{0pt}
\setlength{\parskip}{0.8em}

\pagestyle{fancy}
\fancyhf{}
\cfoot{\thepage}
\renewcommand{\headrulewidth}{0pt}
\renewcommand{\footrulewidth}{0pt}

\newcommand{\code}[1]{\texttt{\small #1}}

\begin{document}

\begin{center}
    \vspace*{5cm}
    {\LARGE \textbf{Michael Polley}} \\[2em]
    \normalsize Student ID: T00774397 \\[2em]
    NFL Big Data Bowl 2027 \\[2em]
    Executive Guide: Beyond the Stopwatch (ELI5 Summary)
    \vfill
\end{center}
\thispagestyle{empty}
\newpage

\noindent\textbf{1.}
\noindent\textbf{(a)}
We consider why traditional scouting at the NFL Scouting Combine fails to identify elite football players. For decades teams have paid millions in draft capital for players who post blazing 40-yard dash times in shorts. We compare this to buying a drag racing car for an obstacle course. A drag racer reaches high speeds in a straight line with nobody around. Real football is an obstacle course where eleven defenders attempt to tackle the ball carrier. The 40-yard dash measures only total elapsed time:
\[ T = \int_{0}^{40} \frac{dx}{v(x)} \]
This single number ignores how an athlete accelerates off the line, changes direction, or handles contact. In our statistical analysis across rookie draft classes, the 40-yard dash time has a $p$-value of $0.343$ when predicting career snaps. Thus straight-line stopwatch speed has zero statistically significant relationship with NFL playing time.

\noindent\textbf{(b)}
We examine where football is actually played on the field. Scouts often look for track athletes who reach top sprint speeds above $32\text{ km/h}$ ($20\text{ mph}$). Across 314,197 regular-season NFL plays in our dataset the median player speed during a play is only $2.06\text{ yd/s}$, which corresponds to $6.78\text{ km/h}$ ($4.2\text{ mph}$). In regular-season competition 95\% of all tracking frames occur below $22.4\text{ km/h}$ ($13.9\text{ mph}$). Speeds exceeding $26\text{ km/h}$ happen on less than 2.5\% of offensive snaps. We define the Core Working Band (CWB):
\[ \text{CWB} \equiv s \in [10.0, 26.0]\text{ km/h} \quad (6\text{ to }16\text{ mph}) \]
Football is played in second and third gear rather than top gear. Elite players like Puka Nacua spend 63.5\% of their Combine drill frames inside this working band. Straight-line sprinters spend their time reaching a high-speed tail that is rarely utilized on Sundays.

\noindent\textbf{2.}
\noindent\textbf{(a)}
We analyze why stopping ability makes an athlete faster than a big engine. In route running a wide receiver must sprint toward a defender, stop abruptly, change direction at a sharp angle, and catch the football. Top speed is useless if an athlete cannot stop. We compare two rookie prospects executing the Combine Comeback drill:
\begin{itemize}[nosep]
    \item Puka Nacua (5th-round draft pick, 1st-Team All-Pro): He exhibits peak braking deceleration of $-5.77\text{ yd/s}^2$. He arrests $708\text{ kg}\cdot\text{m/s}$ of forward momentum in two strides across $0.8\text{ seconds}$ and snaps his hips around to exit the cut at $17.8\text{ km/h}$.
    \item Quentin Johnston (1st-round draft pick, rotational player): He exhibits peak braking of $-5.23\text{ yd/s}^2$. He requires four choppy strides across $1.0\text{ second}$ to arrest his momentum. To avoid losing balance he slows down to $12.4\text{ km/h}$ through the turn.
\end{itemize}
Top performers do not round their routes. They generate greater eccentric braking force to reverse direction while preserving their kinetic energy.

\noindent\textbf{(b)}
We evaluate how velocity preservation predicts career outcomes across 106 Combine wide receivers.

\begin{center}
\small
\setlength{\tabcolsep}{3.5pt}
\begin{tabular}{lcccc}
\toprule
\textbf{Player Tier} & \textbf{Count} & \textbf{Mean Cut Speed ($s_{\text{cut}}$)} & \textbf{Career Status} & \textbf{Outcome} \\
\midrule
Top Performers & 16 & $17.84\text{ km/h}$ & Multi-year Starters & All-Pro / Pro Bowl \\
Average Contributors & 38 & $17.04\text{ km/h}$ & Core Rotation & $3\text{ to }19$ Starts \\
Fringe Prospects & 52 & $16.30\text{ km/h}$ & Backups / Released & $\le 2$ Starts \\
\bottomrule
\end{tabular}
\end{center}

Welch's two-sample $t$-test between Top and Fringe prospects yields $t = 3.723$ and $p = 0.00041$. In our linear regression model each $1.0\text{ km/h}$ of speed preserved through cuts earns an athlete $132$ additional regular-season snaps ($p = 0.002$). In our logistic regression model each $1.0\text{ km/h}$ increase raises the odds of becoming an established NFL starter by 77\% ($p = 0.010, \text{OR} = 1.77$). Thus in-cut velocity preservation is the primary physical prerequisite for NFL wide receivers.

\newpage

\noindent\textbf{3.}
\noindent\textbf{(a)}
We extend our physical framework to defensive linemen and pass rushers using Newton's second law:
\[ \mathbf{F}_{\text{net}} = m \mathbf{a} \implies \mathbf{a} = \frac{\mathbf{F}_{\text{net}}}{m} \]
If a $190\text{ lb}$ athlete accelerates quickly the physical effort is modest because their mass is light. When a $300\text{ lb}$ defensive lineman accelerates quickly off the line of scrimmage the ground reaction force is massive. Because heavier athletes carry greater inertia, raw acceleration naturally decays as body weight increases. To evaluate explosive capacity independent of size, we define Gamma-Over-Expected ($\gamma_{\text{OE}}$). We calculate expected acceleration for a given body weight and subtract it from observed acceleration:
\[ \gamma_{\text{OE}, i} = \gamma_{\text{observed}, i} - \mathbb{E}[\gamma_i \mid \text{Mass}_i] \]

\noindent\textbf{(b)}
We fit a logistic regression model predicting starter status for 114 defensive linemen. The coefficient for body weight alone has $p = 0.614$. The coefficient for $\gamma_{\text{OE}}$ is statistically significant with $p = 0.042$ and odds ratio $\text{OR} = 1.85$. Each $1.0\text{ km/h/s}$ of mass-adjusted acceleration increases the odds of becoming an NFL starter by 85\%. In our tracking data Braden Fiske ($292\text{ lbs}, \gamma_{\text{OE}} = +2.37$) and Laiatu Latu ($259\text{ lbs}, \gamma_{\text{OE}} = +1.70$) earned immediate starting roles due to excess acceleration relative to their mass.

\noindent\textbf{4.}
\noindent\textbf{(a)}
We evaluate defensive backs and offensive linemen to test whether this principle holds across all positions.
\begin{itemize}[nosep]
    \item Defensive Backs ($N=101$): In the 45-degree transition drill, 90th-percentile break speed significantly predicts career defensive snaps ($\beta = +191.84, p = 0.016$). The 40-yard dash is statistically insignificant ($p = 0.471$).
    \item Offensive Linemen ($N=95$): In the pass protection mirror drill, lateral $\gamma_{\text{OE}}$ significantly predicts offensive snaps ($\beta = +154.98, p = 0.026$). The 40-yard dash is statistically insignificant ($p = 0.348$).
\end{itemize}
Across all four major position groups, the 40-yard dash fails while high-frequency kinematic derivatives succeed.

\noindent\textbf{(b)}
We formulate three clear draft-day decision rules for NFL General Managers:
\begin{enumerate}[nosep]
    \item Reject Workout Warriors: When a wide receiver runs a sub-4.40 forty but records an in-cut speed below $16.0\text{ km/h}$, do not draft them on Day 1 or Day 2. Their probability of starting is under 18\%.
    \item Target Heavy Explosion: Prioritize $290+\text{ lb}$ linemen with positive $\gamma_{\text{OE}} \ge +1.50\text{ km/h/s}$. Heavy mass moving with high initial acceleration creates immediate pocket disruption.
    \item Draft Operating Density: Select prospects whose speed density concentrates inside the $10\text{--}26\text{ km/h}$ Core Working Band. Games are decided by intermediate velocity preservation rather than straight-line sprint tails.
\end{enumerate}
Thus evaluating prospects through continuous speed densities and Newtonian mechanics isolates true functional athleticism and eliminates stopwatch draft errors.

\end{document}
